\documentclass[letterpaper]{article}
\usepackage{times}
\usepackage{helvet}
\usepackage{courier}
\usepackage[hyphens]{url}
\usepackage{graphicx}
\usepackage{booktabs}
\usepackage{natbib}
\usepackage{caption}
\urlstyle{rm}
\def\UrlFont{\rm}
\frenchspacing

\title{Constraint Architecture: Post-Hoc Extraction of a Cross-Domain\\
       Agent Safety Skeleton\\
       {\small Submitted to IEEE Software Special Issue --- Building Trustworthy Software in the Time of AI}}

\author{
  Xu Renwu (续仁舞)\\
  Huaqiao University\\
  1nour2567@gmail.com
}

\begin{document}
\maketitle

\begin{abstract}
The 2025--2026 period has produced a clear consensus: LLM-based agents cannot
reliably self-enforce safety rules under adversarial pressure, and multi-layer
deterministic governance is the emerging architectural response.  Multiple
frameworks---LGA, SARC, CAAF, ePCA, ShieldAgent---validate this pattern on
benchmarks.  This paper contributes complementary evidence through a
\textit{post-hoc extraction} approach: two fully functional agent systems, built
by the same developer without shared code in completely unrelated domains
(security operations and embodied AI music), were discovered after completion to
share the same four-layer structure and four structural invariants.  The
invariant was extracted and published as \texttt{constraint-architecture}
(v1.5.1 on PyPI)---a zero-dependency, 1384-line Python skeleton with 7~ABCs +
7~Protocols, 10~frozen dataclasses with \texttt{\_\_post\_init\_\_} value-bound
enforcement, and 25~automated self-checks.  We provide ablation-validated
attack results (16 red-team attacks, 0 breaches; per-layer marginal contribution
estimates), cross-domain structural mapping (security ops, music, finance,
healthcare), and a one-click demo (\texttt{python -m constraint\_architecture.demo}).
We do not claim priority on multi-layer architectures; we claim that convergent
evolution from unrelated practice, combined with an engineering-grade artifact
that enforces architectural commitments at the code level, provides a form of
evidence that benchmark-only validation cannot.
\end{abstract}

% ====================================================================
\section{Introduction}
% ====================================================================

AI agents are moving from research prototypes into production faster than the
safety infrastructure required to constrain them.  Two lines of evidence
establish that prompt-based safety is insufficient.  First, adversarial
testing: Andriushchenko et al.\ \cite{andriushchenko2025} demonstrated 100\%
attack success rates against GPT-4o, Claude~3, and Llama-3 using adaptive
jailbreak attacks.  Second, architectural analysis: the IETF Asimov Safety
Architecture (ASA) draft~\cite{baysal2026} demonstrated that a single LLM
shows progressive relaxation of safety rules under multi-turn pressure---actions
refused at turn~3 were permitted by turn~15, with only accumulated
conversational context changed.  The OWASP LLM01:2025 entry~\cite{owasp2025}
explicitly states that ``it is unclear if there are fool-proof methods of
prevention for prompt injection.''

In 2025--2026, an architectural consensus emerged: \textit{the LLM must not be
the final authority on what executes}.  Multiple independent frameworks
validate the same architectural pattern:

\begin{itemize}
\item LGA \cite{ge2026} --- four-layer governance (Execution Sandbox, Intent
      Verification, Zero-Trust Authorization, Immutable Audit), 96\%
      interception rate across all layers.
\item SARC \cite{besanson2026} --- constraints as first-class specification
      objects, four enforcement sites (Pre-Action Gate, Action-Time Monitor,
      Post-Action Auditor, Escalation Router), zero hard-constraint violations
      under exact predicates.
\item CAAF \cite{li2026} --- ``architecture supersedes model size''; Recursive
      Atomic Decomposition, State Locking (monotonicity), Unified Assertion
      Interface.
\item ePCA \cite{yang2026} --- provably secure guardrails, 0.3~ms verification
      latency, Theorem~1 (Safety Preservation) guarantees no semantic drift.
\item ShieldAgent \cite{zhao2025} --- natural-language safety policies compiled
      to verifiable Python guardrails, 90.4\% accuracy, 3.2$\times$ fewer API
      queries.
\item Five-Plane Architecture \cite{tallam2026} --- stop-anywhere mediation,
      single-digit microsecond adjudication.
\item Three-Layer Assume-Guarantee (Watkins et al.)~\cite{watkins2026} ---
      formal proof that no fewer than three layers is structurally sound.
\end{itemize}

All of these frameworks share a common methodology: \textit{derive principles,
specify layers, build a reference implementation, validate on benchmarks.}
This is rigorous and necessary.  It is not, however, the only form of evidence.

This paper contributes complementary evidence through a different path.
Rather than designing an architecture top-down and validating on benchmarks, we
built two fully functional agent systems in unrelated domains, discovered
\textit{after completion} that they shared the same four-layer structure,
extracted the shared invariant, and published it as a pip-installable
engineering artifact with automated structural self-checks.

Our contribution is \textbf{not} that we discovered multi-layer governance
independently---the frameworks above did it earlier and more rigorously.
Our contribution is threefold:

\begin{enumerate}
\item \textbf{Post-hoc structural convergence.} Two systems---Kylin-Agent
      (security ops) and Malio (embodied music)---converged on the same
      four-layer architecture and four structural invariants without shared
      code or design docs.  We explicitly acknowledge single-author limitation;
      the convergence is suggestive, not dispositive.
\item \textbf{Engineering-grade artifact.} The extracted invariant is not a
      whitepaper---it is \texttt{pip install constraint-architecture}: zero
      dependencies, frozen dataclasses that prevent downstream mutation,
      \texttt{\_\_post\_init\_\_} value-bound enforcement, and 25 automated
      self-checks that run in CI.
\item \textbf{Ablation-validated layering.} Attack-log analysis provides
      per-layer marginal contribution estimates, showing that each layer
      catches attacks the others miss, and that no subset of three layers
      covers all threat categories.
\end{enumerate}

% ====================================================================
\section{The Constraint Architecture}
% ====================================================================

\subsection{Four Deterministic Layers}

The architecture wraps every LLM reasoning step in four deterministic layers.
The principle is simple: \textbf{LLM proposes; code decides; audit proves.}

\begin{description}
\item[Layer 1 -- Reasoning (LLM).] The LLM reasons freely about what to do.
      It cannot execute, authorize, or bypass any downstream layer.  Its output
      is treated as a structured \textit{proposal}, not an instruction.
\item[Layer 2 -- Constraints (Code).] Deterministic validation.  Rules are
      append-only: they may be tightened as trust grows, but never removed.
      Constraints are evaluated against a typed constraint engine.
      Violations raise \texttt{VetoError(reason, alternative, ref)} ---
      structured, machine-consumable, auditable.
\item[Layer 3 -- Execution (Sandbox).] Tiered execution:
      \texttt{AUTO} (read-only diagnostics, immediate),
      \texttt{CONFIRM} (write operations, human approval required),
      \texttt{VETO} (destruction, never executable through this proxy).
      Role and permission enforcement happens here, not in the LLM.
\item[Layer 4 -- Audit (Immutable).] Every pipeline event is recorded in a
      SHA256 hash chain.  \texttt{AuditEvent.prev\_hash} is mandatory at
      construction.  The chain is verifiable at any point; a break is detected
      deterministically.
\end{description}

\subsection{Key Design Invariant: Constraint Monotonicity}

A critical design invariant distinguishes this architecture from permissive
rule systems: \textbf{constraint monotonicity}.  Base safety rules never
decrease.  ``Relaxed'' operational mode means fewer \textit{additional}
restrictions---not weaker baselines.  This is a structural guarantee enforced
by the append-only constraint store, not a prompt-level promise.

This invariant directly addresses the ASA draft's observation that a single
LLM progressively relaxes safety compliance under multi-turn pressure
\cite{baysal2026}.  When the constraint store is append-only and the
validation path is deterministic, no amount of conversational context can
erode the baseline.

\subsection{Frozen Dataclasses as Architectural Commitment}

A subtle but important engineering decision: 10 of the skeleton's 12
dataclasses are \texttt{frozen=True}.  A \texttt{ValidationResult(allowed=False)}
must not be silently flipped to \texttt{allowed=True} by downstream code.
Freezing prevents this at the Python runtime level.  \texttt{\_\_post\_init\_\_}
additionally validates value bounds (e.g., \texttt{risk\_score} bounded to 0--9)
and structural invariants (e.g., \texttt{prev\_hash} is mandatory for
\texttt{AuditEvent}).  These are not documentation promises---they are
runtime-enforced contracts.

% ====================================================================
\section{The Skeleton: Engineering Properties}
% ====================================================================

The extracted invariant is published as \texttt{constraint-architecture} on
PyPI (MIT license).  It is 1384~lines of Python with \textbf{zero
dependencies}.  Table~\ref{tab:engineering} summarizes its structural
properties.

\begin{table}[h]
\centering
\caption{Skeleton v1.5.1 engineering properties.}
\label{tab:engineering}
\begin{tabular}{lrl}
\toprule
Category & Count & Key Detail \\
\midrule
ABCs (abstract classes) & 7 & Shell, Identity, Scheduler, Memory, Protocol, \\
   &    & \quad Tools, + 4 constraint-layer ABCs \\
Protocols (duck-typing) & 7 & One per ABC --- supports inheritance-free composition \\
Frozen dataclasses & 10 of 12 & \verb|__post_init__| validates \verb|risk_score| (0--9), \verb|prev_hash| \\
Custom exceptions & 4 & \verb|VetoError| (reason+alternative+ref), \\
   &    & \quad \verb|ExecutionRefused|, \verb|AuditIntegrity| \\
Enums & 4 & Tier (3 values), EventType (8), MemoryLevel (3) \\
Automated self-checks & 25 & ABC instantiation, frozen immutability, value bounds \\
\addlinespace
Lines of Python & 1384 & Zero dependencies \\
Python version & 3.10+ & \\
Wheel size & 17\,KB & \\
\bottomrule
\end{tabular}
\end{table}

\textbf{ABC + Protocol dual interfaces.}  Every subsystem exposes both an
abstract base class (for inheritance-based wiring) and a Protocol (for
structural subtyping / duck-typing).  Different domains wire the layers
differently: Kylin-Agent uses ABC inheritance throughout; Malio uses
Protocol-based composition.  The skeleton does not force one path.

\textbf{Automated self-checks.}  Running \texttt{python -m
constraint\_architecture.skeleton} executes 25 structural checks: every ABC
must be uninstantiable, every frozen dataclass must reject mutation, value
bounds must be enforced, \texttt{prev\_hash} must be mandatory for
\texttt{AuditEvent}.  If any check fails, the build fails.  This guarantees
that the skeleton itself---not downstream implementations---maintains its
architectural promises.

% ====================================================================
\section{Reference Implementation 1: Kylin-Agent}
% ====================================================================

Kylin-Agent is a security-hardened AI operations agent deployed on Kylin OS
V11---a production-grade Linux distribution widely used in Chinese government
and defense environments.  The agent accepts natural-language operational
commands (e.g., ``check disk usage and restart nginx if it's safe'') and
passes them through four deterministic defense layers:

\begin{description}
\item[T0 -- Injection Detection.] Intercepts and blocks prompt injection
      payloads before the LLM sees them.  Covers direct injection,
      hex/base64/Unicode obfuscation, overflow payloads, and pattern-anomalous
      inputs.  Deterministic---pattern and structural checks only, no LLM
      involved.
\item[T1 -- Risk Scoring.] Assigns a deterministic 1--10 risk score based on
      checklist heuristics.  No ML.  Risk factors include: command
      destructiveness (rm, format, dd), privilege escalation intent,
      target sensitivity (system directories), and anomaly relative to
      historical command patterns.
\item[T2 -- Dual-Path Validation.] Every LLM output is validated through
      \textit{two independent paths}: a structured tool-semantic validator
      (does the LLM's proposed command make semantic sense?) and a raw regex
      validator (does the string match expected command patterns?).  Both
      must agree.  Disagreement triggers human confirmation regardless of
      risk score.  This is the layer that catches linguistic camouflage:
      LLM outputs that \textit{sound safe} but contain dangerous parameter
      combinations.
\item[T3 -- Sandbox.] \texttt{sudo -n} execution with a whitelist.
      Commands not on the allowlist are rejected at the OS level.
      No sandbox bypass is possible without kernel privilege escalation.
\end{description}

After action: T4 audit chain with SHA256 hash linking.  Persistence across
reboots.

\textbf{Red-team results.}  Sixteen independent red-team attacks were conducted.
No attack breached all four layers.  Table~\ref{tab:redteam} (in the
Supplementary Material) provides the full taxonomy.  Key findings:
(1)~T0 intercepted 100\% of obfuscated payloads.  (2)~T2 was critical: in 3 of
16 attacks, the LLM output passed semantic validation but failed raw regex
(dual-path disagreement).  Without the regex path, those 3 would have reached
the sandbox.  (3)~The only audit-chain replay attack was caught at T4.
(4)~No attack reached or breached T3.

% ====================================================================
\section{Reference Implementation 2: Malio}
% ====================================================================

Malio is an embodied AI music agent with 800~particles governed by 9~physics
systems.  Its PersonaEngine enforces a central-bank emotional currency model:
no code can directly modify \texttt{energy}, \texttt{warmth}, or
\texttt{playfulness} parameters.  Agents submit ``loan requests'' for emotional
expressions; the engine grants or denies them deterministically based on
current persona bounds.

\textbf{The same four layers manifest differently:}

\begin{itemize}
\item \textbf{Layer 1 (Reasoning).} MusicAgent (ReAct loop) searches independent
      music; VisualAgent (rule engine) manages the DSL rule lifecycle.
      LLMAutonomous handles proactive heartbeat triggering.
\item \textbf{Layer 2 (Constraints).} PersonaEngine enforces structural bounds:
      e.g., \texttt{light\_burst} (a high-energy visual effect) requires
      energy $\ge$~0.5.  The LLM proposes; the engine vetoes or approves
      deterministically.  In one concrete test, the MusicAgent proposed
      \texttt{light\_burst} while the persona's energy was at 0.2.  The LLM,
      unconstrained, would have executed the action.  The PersonaEngine
      blocked it deterministically.  The LLM did not self-censor.
\item \textbf{Layer 3 (Execution).} Core action push with role filtering.
      No agent can modify another agent's particle parameters directly.
      Visual effects are dispatched through the rendering pipeline, not
      bypassing it.
\item \textbf{Layer 4 (Audit).} Federated rule exchange via SBERT semantic
      clustering.  Every rule proposal has a signed origin and a verifiable
      audit trail.  Conflicting rules from different agents are resolved
      deterministically at the cluster level.
\end{itemize}

\textbf{Why Malio matters.}  Nothing about ``music DJ'' suggests four-layer
security architecture.  The convergence is structural, not semantic.  Both
Kylin-Agent and Malio independently arrived at: (a)~an append-only constraint
store, (b)~tiered execution with a confirm middle tier, (c)~prev\_hash-chained
audit events, and (d)~role enforcement at the code layer.  These four
invariants emerged without a shared codebase, shared design document, or
knowledge of each other's existence during development.

We explicitly acknowledge that single-developer authorship weakens the
independence claim: the same developer's design intuitions may have influenced
both systems unconsciously.  The convergence is suggestive, not dispositive.
A third reference implementation by an independent team would provide
stronger evidence, and the pip-installable skeleton is specifically designed
to minimize the barrier to such external validation.

% ====================================================================
\section{Ablation Analysis and Per-Layer Contribution}
% ====================================================================

A central question for any multi-layer architecture is whether the layers are
redundant.  We provide per-layer marginal contribution estimates based on
attack-log analysis of the 16 red-team attacks against Kylin-Agent
(Table~\ref{tab:redteam} in the Supplementary Material).

\begin{table}[h]
\centering
\caption{Estimated marginal contribution of each defense layer.
         ``Attacks missed'' counts attacks that would reach or breach
         the remaining system if that layer were removed.}
\label{tab:ablation}
\begin{tabular}{lccp{5cm}}
\toprule
Layer Removed & Attacks Missed & Breach Estimate & Attack Types Lost \\
\midrule
None (full 4-layer) & 0/16 & 0\% & --- \\
$-$ T0 (injection detection) & $\ge$4/16 & $\ge$25\% & Direct injection, hex, Unicode, overflow \\
$-$ T2 (dual-path validation) & $\ge$3/16 & $\ge$19\% & Role escalation, tool-parameter, multi-turn \\
$-$ T4 (audit chain) & $\ge$1/16 & $\ge$6\% & Audit-chain replay \\
$-$ T1 (risk scoring) & $\sim$0/16 & $\sim$0\% & (Defense-in-depth; no unique catch in 16 attacks) \\
\bottomrule
\end{tabular}
\end{table}

\textbf{Key finding: no subset of three layers catches all 16 attacks.}
Removing T0 exposes all obfuscated payloads.  Removing T2 allows the 3 attacks
caught by dual-path disagreement to reach the sandbox.  Removing T4 eliminates
replay detection.  T1 provides defense-in-depth---it scored every attack
correctly---but did not uniquely catch any attack in this 16-attack sample;
removing T1 alone might not increase breach rate, but removing T1 \textit{and}
T2 together would compound the failure modes.

These estimates are directionally sound but await formal head-to-head
measurement in a controlled A/B deployment.  They establish that the layers
are not redundant decorators; each catches a distinct class of attacks.

% ====================================================================
\section{Cross-Domain Structural Mapping}
% ====================================================================

If the four-layer architecture is a genuine engineering invariant, it should
map cleanly onto domains \textit{beyond} the two reference implementations.
Table~\ref{tab:crossdomain} demonstrates this with two additional domains.

\begin{table}[h]
\centering
\caption{Four-layer architecture mapped across four domains.}
\label{tab:crossdomain}
\begin{tabular}{lcccc}
\toprule
Domain & Layer 1 (Reasoning) & Layer 2 (Constraints) & Layer 3 (Execution) & Layer 4 (Audit) \\
\midrule
Security ops & LLM diagnoses issues & Block destructive cmds & sudo -n whitelist & SHA256 chain \\
Music (Malio) & MusicAgent searches & PersonaBounds check & Core action push & SBERT-signed rules \\
Finance & LLM selects stocks & Circuit breaker gates & Position size limits & Trade hash chain \\
Healthcare & LLM suggests treatments & Block contraindications & Human confirm require & HIPAA access log \\
\bottomrule
\end{tabular}
\end{table}

In each domain, the same structural invariants hold: the LLM proposes, the
constraint layer gates deterministically, the execution layer enforces
tiered authorization, and the audit layer provides immutable traceability.
The domain-specific content differs; the structural skeleton is invariant.

% ====================================================================
\section{Related Work}
% ====================================================================

\subsection{The Failure of LLM Self-Enforcement}

The ASA draft~\cite{baysal2026} provides the most detailed catalog of
LLM self-enforcement failure modes: frame vulnerability (fictional or
hypothetical frames cause refusal reversal), sycophancy override (helpfulness
pressure overrides security rules), and deterministic pattern denylists with
100\% miss rates on unprogrammed attack patterns.  Andriushchenko et al.\
\cite{andriushchenko2025} confirm that adaptive jailbreak attacks achieve 100\%
success rates against frontier models.  The OWASP LLM01:2025 entry
\cite{owasp2025} concludes that prompt injection has no fool-proof prevention.
Watkins et al.\ \cite{watkins2026} provide a formal result: no genuine
two-stage design can close the safety gap using LLM reasoning alone---a
third structural layer is required.  Together, these results establish that
\textit{LLM self-regulation is structurally insufficient}, motivating the
architectural turn.

\subsection{Multi-Layer Governance Frameworks}

As discussed in Section~\ref{sec:introduction}, LGA, SARC, CAAF, ePCA,
ShieldAgent, and the Five-Plane Architecture have converged on the same
architectural pattern.  We do not claim priority on this pattern; we build
on it, and our contribution is the \textit{post-hoc extraction from practice}
complementing the top-down design from theory.  We note that no existing
framework reports cross-domain structural convergence within a single
codebase---all validate on benchmarks or case studies within a single domain.

\subsection{Layer-Isolated Evaluation}

Hu et al.\ \cite{hu2026} introduce a method for decomposing a production
LLM agent into 8~architectural layers and testing each with its own
deterministic, no-LLM assertion suite (225~cases in 2.39~s, \$0 LLM cost).
Key findings: (1)~aggregate pass-rates mask per-slice failures of 25--91~pp,
(2)~a coverage-honesty criterion refuses to launder uncovered layers,
(3)~the method is a drift detector, not a correctness oracle.  Our skeleton's
25 automated self-checks and frozen-dataclass invariants provide the
structural guarantees that make such layer-isolated evaluation meaningful:
a green slice under layer-isolated evaluation only certifies ``unchanged
from baselined contract''---our skeleton's runtime checks certify that the
contract itself is structurally enforced.

\subsection{Industry Adoption}

Industry frameworks corroborate the architectural consensus.  Microsoft's
FIDES system provides formally verifiable, label-based deterministic defense
for MCP-style agents.  Baidu's four-layer control architecture reports a
failure rate reduction from 18\% to 2.3\% in production deployments
\cite{baidu2026fourlayer}.  The agent-control-plane PyPI package treats the
LLM as ``raw compute'' with a kernel-like governance layer, achieving 0\%
safety violations vs.\ 26.67\% for prompt-based safety.  PolicyLayer/Intercept
provides an open-source deterministic MCP policy engine with stateful counters
and YAML-based rules.

% ====================================================================
\section{The One-Click Demo and Video}
% ====================================================================

\texttt{pip install constraint-architecture} installs the skeleton.
\texttt{python -m constraint\_architecture.demo} runs three commands through
a live constraint engine: a safe \texttt{ls} command (auto-executed), a
destructive \texttt{rm -rf} command (blocked deterministically), and a
privileged \texttt{sudo} command (routed to human confirmation).  The entire
demo runs in under one second with zero configuration.

\texttt{python -m constraint\_architecture.skeleton} executes 25 automated
self-checks: verifying that every ABC is uninstantiable, every frozen
dataclass rejects mutation, \texttt{risk\_score} is bounded 0--9,
\texttt{prev\_hash} is mandatory, and custom exception attributes are
correctly populated.

A 5-minute demo video, source code, PyPI listing, and the full red-team
taxonomy are available at:

\begin{center}
\url{https://pypi.org/project/constraint-architecture/}
\qquad
\url{https://github.com/1nour2567/Constraint-Architecture}
\end{center}

% ====================================================================
\section{Discussion and Limitations}
% ====================================================================

\textbf{Single-author confounding.} The most significant limitation of the
post-hoc convergence argument is single-developer authorship.  The same
developer's design intuitions, coding habits, and safety philosophy may have
influenced both systems unconsciously.  We do not claim that the convergence
proves the architecture's universality---only that it is evidence consistent
with the hypothesis, and stronger than a single-domain case study.

\textbf{Ablation estimates are not controlled experiments.} The per-layer
contribution estimates in Section~\ref{sec:ablation} are based on attack-log
analysis, not a controlled A/B experiment with systematic layer
enable/disable.  A formal ablation study with per-layer toggle and attack
re-execution is planned.

\textbf{Third-domain validation needed.} The cross-domain mapping
(Table~\ref{tab:crossdomain}) is a design exercise, not an implementation.
A third team independently building a constraint-architecture-based system in
finance, healthcare, or education would provide the strongest evidence.

\textbf{Toward standardization.} The convergence of the 2025--2026 literature
toward multi-layer deterministic governance suggests that
\textit{LLM proposes, code decides, audit proves} may be on a trajectory toward
standardization analogous to HTTPS for web security---not because it is
theoretically required, but because practice in security-critical domains
keeps arriving at it independently.  If this trajectory continues, the
\texttt{constraint-architecture} skeleton's role is not to compete with LGA,
SARC, or CAAF---it is to provide a minimal, zero-dependency, copy-paste-able
entry point so that any developer can add a deterministic safety layer to an
existing agent in under five minutes.

% ====================================================================
\section{Conclusion}
% ====================================================================

We presented \texttt{constraint-architecture} (v1.5.1), a cross-domain
validated, pip-installable reference skeleton for AI agent safety.  The
architecture was not designed top-down and validated on benchmarks---it was
\textit{extracted post-hoc} from two independently built, fully functional
agent systems in unrelated domains, published as zero-dependency code with 25
automated self-checks, and validated through 16 red-team attacks (0 breaches)
with per-layer ablation estimates.  We explicitly acknowledge the limitations
of single-author convergence and invite external validation.  Our contribution
is not priority on multi-layer architectures---it is complementary evidence,
arrived at from a different direction, and published as installable code rather
than a policy document.

% ====================================================================
% References
% ====================================================================

\bibliographystyle{plain}
\begin{thebibliography}{}

\bibitem[Andriushchenko et al., 2025]{andriushchenko2025}
Andriushchenko, M.; and Flammarion, N.
Adaptive Jailbreaking of LLMs.
{\em ICLR}, 2025.

\bibitem[Baidu, 2026]{baidu2026fourlayer}
Baidu Developer Center.
Four-Layer Control Architecture for AI Agent Safety.
{\em Baidu Developer Center}, May 2026.

\bibitem[Baysal et al., 2026]{baysal2026}
Baysal, M.; et al.
Asimov Safety Architecture (ASA),
{\em IETF draft-baysal-asimov-safety-architecture-00}, 2026.

\bibitem[Besanson et al., 2026]{besanson2026}
Besanson, G.; et al.
SARC: A Governance-by-Architecture Framework for Agentic AI.
{\em arXiv:2605.07728}, 2026.

\bibitem[Ge et al., 2026]{ge2026}
Ge, J.; et al.
Governance Architecture for Autonomous Agent Systems.
{\em arXiv:2603.07191}, 2026.

\bibitem[Hu et al., 2026]{hu2026}
Hu, D.; et al.
Layer-Isolated Evaluation: Gating the Deterministic Scaffold of a
Production LLM Agent.
{\em arXiv:2606.11686}, 2026.

\bibitem[Li et al., 2026]{li2026}
Li, J.; et al.
CAAF: Harness as an Asset.
{\em arXiv:2604.17025}, 2026.

\bibitem[OWASP, 2025]{owasp2025}
OWASP Foundation.
LLM01:2025---Prompt Injection.  2025.

\bibitem[Santos, 2026]{santos2026}
Santos, R.
Deterministic Guardrails for Enterprise Agents.
{\em Zenodo:19053079}, 2026.

\bibitem[Tallam, 2026]{tallam2026}
Tallam, S.
A Five-Plane Reference Architecture for Runtime Governance of
Production AI Agents.
{\em arXiv:2606.12320}, 2026.

\bibitem[Watkins et al., 2026]{watkins2026}
Watkins, R.; et al.
Position: A Three-Layer Assume-Guarantee Architecture.
{\em arXiv:2605.18672}, 2026.

\bibitem[Yang et al., 2026]{yang2026}
Yang, T.; et al.
Provably Secure Agent Guardrail.
{\em arXiv:2605.29251}, 2026.

\bibitem[Zhao et al., 2025]{zhao2025}
Zhao, S.; et al.
ShieldAgent: Shielding Agents via Verifiable Safety Policy Reasoning.
{\em arXiv:2503.22738}, 2025.

\end{thebibliography}

\end{document}
